\documentclass{article}
\usepackage[english]{babel}
\usepackage{amsthm}
\usepackage{amsmath, amssymb}

% \newtheorem{theorem}{Theorem}[section]
% \newtheorem{lemma}[theorem]{Lemma}
\newtheorem{innercustomthm}{Theorem}
\newenvironment{theorem}[1]{\renewcommand\theinnercustomthm{#1}\innercustomthm}
  {\endinnercustomthm}

\newtheorem{innercustomlem}{Lemma}
\newenvironment{lemma}[1]{\renewcommand\theinnercustomlem{#1}\innercustomlem}
  {\endinnercustomlem}

\newtheorem{innercustomprob}{Problem}
\newenvironment{problem}[1]{\renewcommand\theinnercustomprob{#1}\innercustomprob}
  {\endinnercustomprob}
\begin{document}
\title{MATH 327: Assignment 6}
\author{Grant Yang}
\maketitle

% https://sites.math.washington.edu//~conroy/m327-general/homework/assignment06.pdf
\begin{theorem}{1}
    Let $(a_n)$ be a positive term sequence. \\
    Suppose $\sum\limits_{n=1}^\infty a_n$ converges. \\
    Then $\sum\limits_{n=1}^\infty \frac{\sqrt{a_n}}{n}$ also converges.
\end{theorem}
\begin{proof}
    Let $(a_n)$ be defined as above. \\
    Note that for all $x,y \in \mathbb R$, we have $(x+y)^2 \ge 0$ and so $x^2 + y^2 \ge -2xy$. \\
    Let $n \in \mathbb N$ be any natural, and take $x = -\sqrt{a_n}$ and $y = 1/n$. \\
    Then, we have $a_n + \frac{1}{n^2} \ge 2\frac{\sqrt{a_n}}{n}$. \\
    Since $\sum a_n$ converges and $\sum \frac{1}{n^2}$ is a convergent $p$-series ($p>1$), we know that \\ $\sum\limits_{n=1}^\infty \frac{1}{2}(a_n + \frac{1}{n^2})$ also converges (see lemmas below). \\
    Additionally, for all $n$, \[\left|\frac{\sqrt{a_n}}{n}\right|=\frac{\sqrt{a_n}}{n} \le \frac{1}{2}\left(a_n+\frac{1}{n^2}\right).\]Thus, by the comparison test, $\sum\limits_{n=1}^\infty \frac{\sqrt a_n}{n}$ converges.
\end{proof}

\begin{theorem}{5 from Previous HW}
    $\sum a_n$ converges if and only if $\sum c \cdot a_n$ converges for $c \neq 0$.
\end{theorem}

\begin{lemma}{About Sums}
    If $\sum a_n$ and $\sum b_n$ both converge, then $\sum (a_n+b_n)$ converges.
\end{lemma}
\begin{proof}
    Let $a_n$ and $b_n$ be defined as above, and define the partial sums \[S_k = \sum\limits_{n=1}^k (a_n+b_n), \quad A_k = \sum\limits_{n=1}^k a_n, \quad B_k = \sum\limits_{n=1}^k b_n.\] \\
    Note that $S_k = A_k + B_k$, and since $\sum a_n$ and $\sum b_n$ converge, $\lim A_k$ and $\lim B_k$ exist. \\
    Then, by the properties of limits, $\lim S_k$ exists and $\lim S_k = \lim A_k + \lim B_k$.
\end{proof}

\newpage{}
\begin{theorem}{2}
    Let $(a_n)$ be a decreasing sequence of positive real numbers. \\
    Suppose $\sum\limits_{n=1}^\infty a_n$ converges. Then $\lim n a_n = 0$. 
\end{theorem}
\begin{proof}
    Let $(a_n)$ be defined as above. \\
    Let $\varepsilon > 0$ be any real number. \\
    By the Cauchy criterion of convergence, there must exist $N \in \mathbb N$ such that for all $m, k > N$ and $k \ge m$, we have
    $\left|\sum\limits_{n=m}^k a_n \right| < \varepsilon/2$. \\
    For any $m>N$, choose $k = 2m$, and since $a_n$ is positive and decreasing,
    \begin{align*}m\cdot a_k\le (m+1)a_k=(k-m+1)\cdot a_k\le \sum_{n=m}^k a_n =\left|\sum_{n=m}^k a_n \right| < \varepsilon/2. \end{align*}
    Thus, $2m\cdot a_k < \varepsilon$, and $2m \cdot a_{2m} < \varepsilon$. \\
    We have shown that the even term subsequence converges with $\lim 2m\cdot a_{2m} = 0$. \\
    Let $n \in \mathbb N$, and now consider the odd term subsequence (starting from $a_3$). \\
    Note that $(2n+1) \cdot a_{2n+1}\le (2n+1)\cdot a_{2n} = 2n\cdot a_{2n} + a_{2n}$. \\
    Since $\lim a_n = 0$ as $\sum a_n$ converges, we have $\lim (2n\cdot a_{2n} + a_{2n}) = 0$. \\
    Thus, $(2n+1)\cdot a_{2n+1}$ is bounded below by $0$ and above by a sequence converging to $0$, so by the squeeze theorem it converges to $0$. \\
    Both the odd term and even term subsequences of $n\cdot a_n$ converge to zero, so $\lim n\cdot a_n = 0$.\footnote{Every term except $a_1$ is either in the odd subsequence or even subsequence, so for every $\varepsilon > 0$ we can just choose the maximum of the two adjusted $N$s.}
\end{proof}

\newpage{}
\begin{problem}{3}
    Determining convergence or divergence:
\end{problem}
\begin{enumerate}
    \item[(a)] $\sum\limits_{n=1}^\infty \frac{n}{2^n}$: Converges by the ratio test. $\lim \frac{n+1}{2^{n+1}} \cdot \frac{2^n}{n} = \frac{1}{2}\lim (1+\frac{1}{n}) = \frac{1}{2}<1$.
    \item[(b)] $\sum\limits_{n=2}^\infty \frac{1}{(\log n)^{\log n}}$: By Cauchy condensation, this converges if and only if the following converges\footnote{Note that we are free to shift around the starting points of the series to make the Cauchy condensation line up where we want. We can define $b_n = a_{n+m}$ for $m \in \mathbb Z$ (defining added terms arbitrarily), and $\sum b_n$ converges iff $\sum a_n$ converges. Then we take Cauchy condensation $2^n b_{2^n}$ and we are free to chop off some terms here too.}: \[\sum\limits_{n=2}^\infty \frac{2^n}{(n \log 2)^{n \log 2}} = \sum_{n=2}^\infty \left(\frac{2}{(n \log 2)^{\log 2}}\right)^n.\]
    By the root test, this series converges since \[\lim \frac{2}{(n \log 2)^{\log 2}} = \frac{2}{(\log 2)^{\log 2}}\lim\frac{1}{n^{\log 2}} = 0 <1.\]
    Since $\log 2>0$ the denominator is an unbounded increasing function, so the limit is $0$.
    \item[(c)] $\sum\limits_{n=1}^\infty (-1)^{n+1} \frac{n}{n+1}$: Diverges since $\lim a_n \neq 0$. In fact that limit does not converge! Note that for all $n>1$, we have $2n> n+1$ and so $\frac{n}{n+1} > \frac{1}{2}$. This means $|a_n| > \frac{1}{2}$ for all $n > 1$, so $\lim a_n \neq 0$.
    \item[(d)] $\sum\limits_{n=1}^\infty \frac{n!}{n^n}$: Converges by the ratio test: \begin{align*} 
    \lim \frac{(n+1)!}{(n+1)^{n+1}}\cdot \frac{n^n}{n!} &= \lim \frac{n+1}{n+1}\left(\frac{n}{n+1}\right)^n \\ &= \lim \frac{1}{\left(1+\frac{1}{n}\right)^n} = \frac{1}{e} < 1.
    \end{align*}
    We're just gonna pretend we know this limit is $1/e$, trust. 
    \item[(e)] $\sum\limits_{n=1}^\infty \frac{(n!)^2}{(2n)!}$: Converges by the ratio test: \[\lim\frac{[(n+1)!]^2}{(2n+2)!}\cdot \frac{(2n)!}{(n!)^2} = \lim\frac{(n+1)^2}{(2n+2)(2n+1)} = \frac{1}{4}<1.\] \\
    \vspace{1em} \\ 
    Sorry, part (f) is orphaned on the next page.
    \newpage{}
    \item[(f)] $\sum\limits_{n=5}^\infty (-1)^{n+1}\frac{\sqrt n}{n-4}$: Converges by the alternating series test.  \\ For $n > 8$, note that $2(n-4) > n$ and so $\frac{\sqrt{n}}{n-4}< \frac{\sqrt{n}}{n/2} = \frac{2}{\sqrt{n}}$. \\ $\lim \frac{2}{\sqrt{n}} = 0$ as the denominator is an unbounded increasing function. \\ By the squeeze theorem on $0 < \frac{\sqrt{n}}{n-4} < \frac{2}{\sqrt{n}}$ for $n>8$, we get $\lim a_n = 0$. \\ 
    Now, we can take a derivative to show that $\sqrt{n}/(n-4)$ is eventually decreasing, but unfortunately we don't know what a derivative is. \\
    Take that $n>4$. 
    Bashing the inequality, where each subsequent line implies the previous (read it backwards):
    \begin{align*}
        \frac{\sqrt{n+1}}{n+1-4} &\le \frac{\sqrt n}{n-4} \\
        (n-4) \sqrt{n+1}&\le (n-3)\sqrt {n} \\
        (n^2-8n+16)(n+1) &\le (n^2-6n+9)n \\
        16&\le n+n^2.
    \end{align*}
    
    % We don't know what a derivative is, so it seems my only option here is to bash out the inequality?? which seems awful. 
\end{enumerate}


\begin{theorem}{4}
    Let $(a_n)$ be a sequence of reals and suppose that $\sum |a_n|$ converges. \\
    Let $(b_n)$ be a bounded sequence of reals. \\
    Then $\sum a_n b_n$ converges.
\end{theorem}
\begin{proof}
    Let $(a_n)$ and $(b_n)$ be defined as above. \\
    Since $b_n$ is bounded, $|b_n|$ is also bounded. \\
    Let $B$ be the upper bound of $|b_n|$. \\
    Note that if $B = 0$, this means that $|b_n| = 0$ and thus $b_n = 0$ for all $n$. \\
    Then, $\sum\limits_{n=1}^\infty 0\cdot a_n = 0$ as desired, so we only need to consider the remaining case that $B> 0$.\\
    Let $\varepsilon > 0$ be any real. \\
    For any $m,k \in \mathbb N$, we can calculate
    \[\left|\sum_{n=m}^k a_n b_n\right| \le \sum_{n=m}^k |a_n||b_n| \le B\sum_{n=m}^k |a_n|.\]
    By the Cauchy criterion for the convergence of $|a_n|$, let $N \in \mathbb N$ be such that for all $k,m \in \mathbb N$ where $k \ge m > N$, we have $\sum\limits_{n=m}^k |a_n| < \varepsilon/B$. \\
    Then, for any $m,k \in \mathbb N$ with $k \ge m > N$,
    \[\left|\sum_{n=m}^k a_n b_n\right| \le B \sum_{n=m}^k |a_n| < \varepsilon,\]
    so by the Cauchy criterion, $\sum a_n b_n$ converges.
\end{proof}

\newpage{}
\begin{theorem}{5}
    Let $\sum a_n$ and $\sum b_n$ be convergent series of non-negative numbers. \\
    Then $\sum \sqrt{a_n b_n}$ converges.
\end{theorem}
\begin{proof}
    Let $(a_n)$ and $(b_n)$ be defined as above. \\
    First, we can show that for all $x,y\ge 0$, we have $\sqrt{xy} \le x + y$. \\
    Note that since $xy\ge0$ we have $2xy\ge xy$. \\
    Then, since $x^2 +y^2 \ge 0$ we have $(x+y)^2=x^2+y^2+2xy \ge 2xy \ge xy$. \\
    Since both $(x+y)^2$ and $xy$ are positive and the square root is an increasing function, $x+y\ge \sqrt{xy}$ as desired. \\
    Let $n \in \mathbb N$, and take $x = a_n$ and $y = b_n$. \\
    Then, $|\sqrt{a_n b_n}| = \sqrt{a_n b_n} \le a_n + b_n$. \\
    Because both $a_n$ and $b_n$ converge, by the Lemma About Sums proven for Theorem 1 above, $\sum\limits_{n=1}^\infty (a_n+b_n)$ converges. \\ By the comparison test, $\sum\limits_{n=1}^\infty \sqrt{a_n b_n}$ also converges.
\end{proof}

\newpage{}
\begin{theorem}{6}
    Let $\sum\limits_{n=1}^\infty a_n$ be a conditionally convergent series (i.e. $\sum a_n$ converges and $\sum |a_n|$ diverges). \\
    The series made of positive terms of $(a_n)$ and the series made of negative terms of $(a_n)$ each diverge.
\end{theorem}
\begin{proof}
    Let $(a_n)$ be defined as above.  \\
    We wish to show that the series made of negative terms of $(a_n)$ diverges. \\
    First, note that the number of negative terms appearing in $(a_n)$ is infinite: \\
    If there were some last negative term $a_p$, then $\sum\limits_{n=p+1}^\infty a_n = \sum\limits_{n=p+1}^\infty |a_n|$, which means $\sum a_n$ and $\sum |a_n|$ either both converge or both diverge, contradicting the fact that $\sum a_n$ is conditionally convergent. \\
    Define the partial sums $S_k = \sum\limits_{n=1}^k a_n$ and $A_k = \sum\limits_{n=1}^k |a_n|$. \\
    Define $T_k$ as the partial sum of the first $k$ negative terms of $a_n$. \\
    Define $\varphi$ to be a function from $\mathbb N$ to $\mathbb N$, mapping $n$ to the number of negative terms appearing in $a_1,\dots,a_n$. \\
    Note that $\varphi$ is increasing and surjective ($T_k$ is infinite). \\
    Observe that \[A_k = S_k -2T_{\varphi(k)}.\]
    Assume FTSOC that $T_k$ converges, i.e. $\lim T_k = T$ exists. \\
    By the fun lemma below, $\lim T_{\varphi(k)} = \lim T_k = T$. \\
    Since $\sum a_n$ converges, $\lim S_k = S$ exists. \\
    Thus, by the properties of limits, $\lim A_k = S - 2T$ also exists. \\
    However, this contradicts the fact that $\sum |a_n|$ diverges, so our assumption was false and $\lim T_k$ does not exist, and the series of negative terms of $(a_n)$ diverges. \\
    Now, note that $\sum (-1)\cdot a_n$ is also conditionally convergent since $\sum |-a_n| = \sum |a_n|$, and the negative terms of that sequence are exactly the positive terms of $\sum a_n$. \\
    Since the series of the negative terms of $\sum (-1)\cdot a_n$ diverges, we know the series of the positive terms of $\sum a_n$ diverges. \\
    Thus, the series of positive terms of $(a_n)$ and the series of negative terms of $(a_n)$ both diverge.
\end{proof}

\begin{lemma}{(Fun Lemma)}
    Let $(T_k)$ be a sequence, and let $\varphi$ be an increasing surjective function from $\mathbb N$ to $\mathbb N$. \\
    If $\lim T_k = T$, then $\lim T_{\varphi(k)} = T$.
\end{lemma}
\begin{proof}
    Let $(T_k)$, $T$, and $\varphi$ be defined as above. \\
    Let $\varepsilon > 0$, and by the convergence of $(T_k)$ let $N \in \mathbb N$ be such that for all $n > N$, we have $|T_n - T| < \varepsilon$. \\
    Since $\varphi$ is surjective, let $K$ be such that $\varphi(K)=N+1$. \\
    For any $k > K$, we have $\varphi(k) \ge \varphi(K)=N+1>N$ since $\varphi$ is increasing. \\
    Thus, for all $k > K$, $|T_{\varphi(k)} - T| < \varepsilon$.
\end{proof}
\end{document}